\usepackage[utf8]{inputenc}
\usepackage[a4paper,top=2cm, bottom=2cm, outer=7cm, inner=1cm, heightrounded, marginparwidth=5cm, marginparsep=1cm]{geometry}
\usepackage{graphicx}
\usepackage[german]{babel}
%\usepackage{blindtext}
\usepackage{marginnote}
\usepackage{multirow}
\usepackage{diagbox}
\usepackage{mhchem}
\usepackage{url}
\usepackage{bbm}
\usepackage{biblatex}
\addbibresource{sources.bib}
\usepackage{amsfonts,latexsym} % para tener disponibilidad de diversos simbolos
\usepackage{enumerate}
%\usepackage{booktabs}
\usepackage{float}
\usepackage{subcaption}
\usepackage{tikz}
%\usepackage{threeparttable}
\usepackage{array,colortbl}
%\usepackage{ifpdf}
\usepackage{longdivision}
\usepackage{polynom}
\usepackage{xcolor}
%\usepackage{rotating}
%\usepackage{cite}
\usepackage{stfloats}
\usepackage{listings}
\usepackage{fancyhdr}
\usepackage{makeidx}
\usepackage{hyperref}
\hypersetup{
    colorlinks=true,
    linktoc=all,
    linkcolor=CarnationPink,   % Farbe für interne Links
    citecolor=CarnationPink,    % Farbe für Zitierungen
    urlcolor=CarnationPink       % Farbe für URLs
}
%\usepackage{wrapfig}
\usepackage{array}
\usepackage[symbol]{footmisc}
%\usepackage{multirow}
%\usepackage{adjustbox}
%\usepackage{nccmath}
% BOXES 
\usepackage{tikz}
\usepackage[most]{tcolorbox}
\tcbset{blue box/.style={
    enhanced, fonttitle=\bfseries,
    colback=white, colframe=accdarklight,
    coltitle=white, colbacktitle=accdark,
    attach boxed title to top left={xshift=0.3cm,
                                    yshift*=-\tcboxedtitleheight/2},
    boxed title style={
      before upper=\hspace*{0.5cm}, % reserve space for the image
      overlay={
       \node at ([xshift=0.5cm]frame.west) {};
        
      }
    }
  }
}
\newtcolorbox{mybox}[1][]{blue box, #1}
\usepackage{amsmath}
\usepackage{amssymb}
\usepackage[dvipsnames]{xcolor}


\newcommand{\titulo}{Aufarbeitung Diskrete Mathematik}
\newcommand{\nombre}{Lil'Jana}
\newcommand{\FCE}{\textbf{WiSe 2024/2025}}
\newcommand{\resumen}{Abstrakt}


% Glossary 
\usepackage[automake]{glossaries}
\newglossary*{dinge}{Was ist was?}
\glsaddkey
 {formB}
 {}
 {\glsentryB}
 {\GlsentryB}
 {\glsB}
 {\GlsB}
 {\GLSB}
 \glsaddkey
 {formF}
 {}
 {\glsentryF}
 {\GlsentryF}
 {\glsF}
 {\GlsF}
 {\GLSF}
\makeglossaries
\loadglsentries{gloss} 

\newcommand{\charfunc}[1]{\glsentrytext{charfunc} mit $A = #1$, also $\chi_{#1}$}


% Index
\usepackage{imakeidx}
\makeindex




\usepackage{fancyhdr}
\pagestyle{fancy}
\fancyhead{}
\fancyfoot{}
\fancyhead[LO]{\small\nombre}
\fancyhead[RE]{\small\titulo}
\fancyhead[LO]{\small\FCE}
\fancyfoot[RO,LE] {\thepage}

\setcounter{page}{1} %Número de la página, la editorial lo cambiará


\newcommand{\mapsfrom}{\mathrel{\style{display:inline-block; transform:scale(-1,1);}{\mapsto}}}
\newcommand{\X}{$X$}
\usepackage{amsthm}
%\theoremstyle{definition}
\newtheoremstyle{mytheoremstyle} % name
    {\topsep}                    % Space above
    {\topsep}                    % Space below
    {}                   % Body font
    {}                           % Indent amount
    {\bfseries}                   % Theorem head font
    {}                          % Punctuation after theorem head
    {1.5em}                       % Space after theorem head
    {\thmname{#1}\thmnumber{ #2}\thmnote{ #3}}  % Theorem head spec (can be left empty, meaning ‘normal’)

\theoremstyle{mytheoremstyle}
\newtheorem{bsp}{Beispiel}[section]
\newtheorem{info}[bsp]{Zusatzinformation}
\newtheorem{reg}[bsp]{Regel}
\newtheorem{satz}[bsp]{Satz}
\newtheorem{lem}[bsp]{Lemma}
\newtheorem{defi}[bsp]{Definition}
\newtheorem{fr}[bsp]{Frage}
\newtheorem{fakt}[bsp]{Fakt}
\newtheorem{beo}[bsp]{Beobachtung}
\newtheorem{prinz}[bsp]{Prinzip}
\newtheorem{bem}[bsp]{Bemerkung}
\newtheorem{anmerkung}{Anmerkung}[]
\newtheorem{folg}[bsp]{Folgerung}
\newtheorem{er}[bsp]{Erinnerung}
\newtheorem{prob}[bsp]{Problem}
\newtheorem{prop}[bsp]{Proposition}
\newtheorem{beob}[bsp]{Beobachtung}
\newtheorem{eig}[bsp]{Eigenschaften}
\newtheorem{kor}[bsp]{Korollar}
\newtheorem{beh}[bsp]{Behauptung}
\newcommand{\anm}[2]{
\color{Blue} \noindent
\begin{anmerkung}  \textit{zu "#1"}  \newline
#2
\end{anmerkung} \color{black}
}

\newcommand{\regel}[2]{
\begin{reg} \textbf{#1}  \newline
#2
\end{reg}
}
\newcommand{\satz}[2]{
\begin{satz} \textbf{#1}  \newline
#2
\end{satz}
}
\newcommand{\prin}[2]{
\begin{prinz} \textbf{#1}  \newline
#2
\end{prinz}
}
\newcommand{\rsp}[3]{
\noindent \textbf{#1} \pageref{rsp:#2}  

\noindent #3 

}
\newcommand{\txc}[1]{
\textcolor{Blue}{#1}
}
\newcommand{\txca}[1]{
\textcolor{Aquamarine}{#1}
}

\newcommand{\za}[1]{\mathbb{Z} / #1 \mathbb{Z}}
\newcommand{\zat}[0]{\mathbb{Z} / m \mathbb{Z}}
\newcommand{\zatx}[1]{\mathbb{Z} / #1 \mathbb{Z}[x]}
\usepackage[onehalfspacing]{setspace} 
\renewcommand{\thefootnote}{\roman{footnote}}


\newcommand{\kl}[1]{\{ #1 \}}
\newcommand{\sumEN}[0]{\sum_{i=1}^n}
\newcommand{\un}[1]{\underline{#1}}
\newcommand{\nn}[0]{n \times n}
\newcommand{\ma}[1]{\(#1\)}
\newcommand{\mab}[1]{\[#1\]}
\newcommand{\beg}[1]{\un{#1} \index{#1}}
\newcommand{\mnnk}[0]{\mathbb{M}_{\nn}(K)}
\newcommand{\R}[0]{\mathbb{R}}
\newcommand{\Tau}[0]{\mathcal{T}}
\newcommand{\itr}[0]{I \rightarrow \R}
\newcommand{\intab}[0]{\int_a^b}
\newcommand{\isr}[0]{I \subseteq \R}
\newcommand{\dx}[0]{\,dx}
\newcommand{\dy}[0]{\,dy}
\newcommand{\aub}[2]{\txca{
     \underbrace{
     \textcolor{black}{#1}
     }_{#2}
     }}
\newcommand{\bub}[2]{\txc{
\underbrace{
\textcolor{black}{#1}
}_{#2}
}}
\newcommand{\aub}[2]{\txca{
\overbrace{
\textcolor{black}{#1}
}^{#2}
}}
\newcommand{\bob}[2]{\txc{
\overbrace{
\textcolor{black}{#1}
}^{#2}
}}

\usepackage{todonotes}
\newcommand{\bluenote}[1]{\todo[backgroundcolor=white, bordercolor=Blue, color=Blue!10]{\textcolor{Blue}{#1}}}
\newcommand{\aquanote}[1]{\todo[backgroundcolor=white, bordercolor=Aquamarine, color=Aquamarine!10]{\textcolor{Aquamarine}{#1}}}